% 2026-10-21 统一复核：几何、箭头、符号与正文逐项对应。
% 印宽17cm；最小字号8.4pt；保留独立注释行与框间空隙。
% 图契约：本题的源、边界和推导逐项对应；箭头符号按图注解释。
\documentclass[tikz,border=7pt]{standalone}
\usepackage{fontspec,xeCJK,amsmath,amssymb}
\IfFontExistsTF{Times New Roman}{\setmainfont{Times New Roman}}{\setmainfont{TeX Gyre Termes}}
\IfFontExistsTF{Heiti SC}{\setCJKmainfont[AutoFakeBold=1.5]{Heiti SC}}{\setCJKmainfont[AutoFakeBold=1.5]{FandolHei-Regular}}
\usetikzlibrary{arrows.meta,bending,calc,patterns,positioning}
\definecolor{ink}{HTML}{243745}
\definecolor{blue}{HTML}{275DAB}
\definecolor{green}{HTML}{177C63}
\definecolor{red}{HTML}{BD4935}
\definecolor{orange}{HTML}{AA6C16}
\definecolor{purple}{HTML}{715393}
\tikzset{every node/.style={font=\fontsize{9.3}{12.5}\selectfont,text=ink},
axis/.style={-{Stealth[length=4pt]},line width=.65pt,draw=ink!65},
vec/.style={-{Stealth[length=5pt]},line width=1pt,draw=blue},
short/.style={-{Stealth[length=3pt]},line width=.8pt},
guide/.style={dashed,draw=ink!40,line width=.5pt},
card/.style={rounded corners=2pt,draw=blue!65,fill=blue!4,align=center,text width=6.5cm,inner xsep=9pt,inner ysep=7pt},
title/.style={font=\fontsize{10.5}{14}\selectfont\bfseries,text width=15cm,align=center},
label/.style={font=\fontsize{8.4}{11}\selectfont,align=center},
note/.style={label,text width=6.2cm}}
\begin{document}\begin{tikzpicture}[x=1cm,y=1cm]
\path[use as bounding box] (-8,-3.5) rectangle(8,3.5);

\node[title] at(0,3.2){三个时间尺度：先核对响应模型，再比较数值};
\node[card,text width=4.25cm,minimum height=3.5cm] at(-5.25,.6){
\textbf{动量响应}\\[5pt]$\tau_m$：Drude 载流子惯性\\[5pt]
$\tau_m\dot{\mathbf J}+\mathbf J=\sigma_0\mathbf E$\\[5pt]
铜约 $2.5\times10^{-14}\,\mathrm s$};
\node[card,text width=4.25cm,minimum height=3.5cm,draw=green!70,fill=green!4] at(0,.6){
\textbf{局域体电荷衰减}\\[5pt]瞬时欧姆近似：$\tau_c=\varepsilon/\sigma_0$\\[5pt]
一阶衰减须 $\tau_c\gg\tau_m$\\[5pt]
铜的 $1.5\times10^{-19}\,\mathrm s$\\是直流外推值，不是实际短时常数};
\node[card,text width=4.25cm,minimum height=3.5cm,draw=orange!70,fill=orange!4] at(5.25,.6){
\textbf{整体电路建立}\\[5pt]
传播：由长度与波速决定\\[5pt]
集中参数：$RC$、$L/R$\\[5pt]
局域响应不能代替整条线路的建立};
\node[align=center,text width=15cm] at(0,-2.35){保留载流子惯性时：$\displaystyle
\tau_m\ddot\rho_f+\dot\rho_f+\frac{\sigma_0}{\varepsilon}\rho_f=0$\\[4pt]
只有慢变、过阻尼且惯性可忽略时，才退化为 $\rho_f(t)=\rho_f(0)e^{-t/\tau_c}$。};
\end{tikzpicture}\end{document}